\subsection{Data Limitations}
Before proceeding with the analysis, it is important to note some important limitations of the data at hand.  First, the most salient limitation of the data source that I rely on in this study concerns the limited range of cultural domains included in the NSS for which information on the rate of culture-based talk is available.  In particular, as befits an arts participation survey, questions about culture-based talk are weighted toward the most legitimate forms of aesthetic practice but are silent when it comes to more popular culture consumption experiences \citep{dimaggio_mukhtar04}. In particular, past research suggests that culture-based talk centered around sports \citep{erickson96}; Fashion and Cuisine \citep{dimaggio87} and in particular television \citep{fiske87} are a consistent generator of culture-based talk in formal  and informal settings (e.g. the workplace).  Thus, we can be certain that the rates of culture based talk reported below \emph{underestimate} by a substantial amount its prevalence in the American population.   The absence of information on culture-based talk in more popular domains may also lead to an \emph{overstatement} of its connection to markers of stratification, such as educational credentials; although it is reasonable to suppose that we may be likely to observe education-based stratification in rates of culture-based talk in the absence of education-based stratification in consumption rates.  

Another important issue concerns the fact that there is no complete correspondence between the culture-talk and culture-consumption questions.  This becomes relevant whenever we attempt to compare the relative salience of socio-demographic differences across the talk and consumption domains.  This problem is most salient in the case of the ``live performance'' items.  While the culture-talk item deals with the more restricted domain of live drama, the consumption question deals with the broader domain of all live performances, including musicals and dance.  This obviously introduces error into any attempt to compare the relative importance of status markers across the two domains.  In the case of Movies and the Arts, the consumption questions are more restrictive than the talk questions, since the former deal explicitly with engagement with film and art outside of the home (going to the movies, visiting a gallery or museum), while the latter deal with the use of film or arts culture in general as a conversational topic.  It is of course likely that many respondents engage film and the arts without having to go the movie theater and the gallery, and this cultural knowledge should still be available for interaction all the same.  While we should expect broad positive correlations between more restricted and more generalized forms of participation \citep{dimaggio_useem78a, sullivan_katzgerro07}, the lack of strict congruence in the live performance domain should be kept in mind when interpreting the results. 

One final limitation of the data source is the fact that we have to rely on a single set of survey items to measure culture-based talk.  Survey questions have the advantage of allowing us to draw more general conclusions and subject the data to a powerful set of statistical manipulations.  However, they have the disadvantage of glossing the \emph{context} of consumption as a mediator of the extent to which given aesthetic consumption experiences will generate culture-based talk. For instance, some genres have built-in rituals of audience reception in which ``interaction'' is expected to occur \emph{within} the consumption site and \emph{during} the performance (e.g. a rock concert) and which thus carry little expectations that the cultural-knowledge gained therein will be transposed to other informal contexts as a generator of conversation.  For instance, culture that is consumed according to socially constructed standards of audience ``propriety'' (e.g. sitting quietly in rapt attention) on the other hand, may in turn generate a lot of culture-based talk precisely because direct interaction among audience members (verbal and non-verbal) is frowned upon during the reception process.\footnote{Cultural domains that exhibit wide variation in the context (and thus the expected audience conventions) of reception, may be seen to have a ``bi-modal'' distribution of their likelihood of being used for culture-based talk, with non-negligible proportions of active consumers reporting not using (as well as using) culture for that domain as a conversational resource. }  